% Cajas de contenido corporativas (nota, aviso, importante) y estilo del
% resumen. Se inyecta en la cabecera del documento con `pandoc -H` y se
% activan con el filtro pandoc-latex-environment a partir de divs con clase
% (::: nota, ::: aviso, ::: importante). Mismos colores que
% fpvirtual/customization.latex. Las definiciones se protegen con guards para
% no chocar con entornos definidos por el propio documento.
\usepackage[most]{tcolorbox}

\definecolor{fpvbox-orange}{HTML}{F3970D}
\definecolor{fpvbox-darkorange}{HTML}{C77800}
\definecolor{fpvbox-red}{HTML}{A50000}
\definecolor{fpvbox-midgray}{HTML}{808080}

% resumen — caja del abstract (YAML), naranja corporativo
\tcolorboxenvironment{abstract}{
  enhanced, breakable,
  colback=fpvbox-orange!6!white, colframe=fpvbox-orange,
  boxrule=0.5pt, arc=1.2mm,
  left=8pt, right=8pt, top=6pt, bottom=6pt,
  title=Resumen, colbacktitle=fpvbox-orange, coltitle=white,
  fonttitle=\sffamily\bfseries\small, titlerule=0mm,
}

% nota — información complementaria
\ifcsname nota\endcsname\else
\newtcolorbox{nota}{
  enhanced, breakable,
  colback=fpvbox-midgray!7!white, colframe=fpvbox-midgray!40!white,
  boxrule=0.5pt, arc=1.2mm,
  left=8pt, right=8pt, top=6pt, bottom=6pt,
  title=Nota, colbacktitle=fpvbox-midgray!40!white, coltitle=white,
  fonttitle=\sffamily\bfseries\small, titlerule=0mm,
}
\fi

% aviso — precaución
\ifcsname aviso\endcsname\else
\newtcolorbox{aviso}{
  enhanced, breakable,
  colback=fpvbox-darkorange!7!white, colframe=fpvbox-darkorange,
  boxrule=0.5pt, arc=1.2mm,
  left=8pt, right=8pt, top=6pt, bottom=6pt,
  title=Aviso, colbacktitle=fpvbox-darkorange, coltitle=white,
  fonttitle=\sffamily\bfseries\small, titlerule=0mm,
}
\fi

% importante — advertencia grave / irreversible
\ifcsname importante\endcsname\else
\newtcolorbox{importante}{
  enhanced, breakable,
  colback=fpvbox-red!5!white, colframe=fpvbox-red,
  boxrule=0.5pt, arc=1.2mm,
  left=8pt, right=8pt, top=6pt, bottom=6pt,
  title=Importante, colbacktitle=fpvbox-red, coltitle=white,
  fonttitle=\sffamily\bfseries\small, titlerule=0mm,
}
\fi

% ---------------------------------------------------------------------------
% Bloques de código con caja (aspecto Eisvogel) SIN listings: se conserva el
% resaltado de Pandoc (--syntax-highlighting=default) y los caracteres Unicode
% (listings descoloca los multibyte: diagramas con ─ │ └ ▶).
% ---------------------------------------------------------------------------
\usepackage[most]{tcolorbox}
\usepackage{etoolbox}
\definecolor{codigo-fondo}{HTML}{F7F7F7}
\definecolor{codigo-borde}{HTML}{B3B2B3}
\tcbset{codigo/.style={breakable, enhanced, colback=codigo-fondo,
  colframe=codigo-borde, boxrule=0.4pt, arc=0pt, outer arc=0pt,
  left=5pt, right=5pt, top=3pt, bottom=3pt, before skip=0.8em, after skip=0.8em}}
% Bloques con lenguaje resaltado (Shaded/Highlighting de Pandoc).
\AtBeginDocument{\renewenvironment{Shaded}{\begin{tcolorbox}[codigo]}{\end{tcolorbox}}}
% Bloques sin lenguaje o no resaltables (verbatim): misma caja y cortes de línea.
\RecustomVerbatimEnvironment{verbatim}{Verbatim}{breaklines,fontsize=\small}
\BeforeBeginEnvironment{verbatim}{\begin{tcolorbox}[codigo]}
\AfterEndEnvironment{verbatim}{\end{tcolorbox}}
